\section{Reproduction and Artifact Map}\label{sec:reproduction}

\subsection{Three distinct reproduction levels}

A clone contains code, synthetic/golden inputs, selected derived artifacts,
and this report's generation inputs. It does not contain the raw market
capture. The following levels must be distinguished:

\begin{enumerate}
\item \textbf{Mechanics and native implementation.} Install the supported
Python environment, run deterministic tests, build C++, and compare native
behavior on committed/synthetic traces. These do not require raw data.
\item \textbf{Committed evidence.} Verify the historical V2 structure and
current V3 durable artifacts, reconstruct the decision summary from committed
outputs, and regenerate this report. Artifact-only mode verifies recorded
input identities but does not claim to read missing capture bytes.
\item \textbf{Full empirical replay.} Restore all selected raw files at their
manifest paths with matching SHA-256 values, use the frozen research source,
and rerun the panel. Native pipeline comparisons similarly require selected
raw development captures.
\end{enumerate}

The project is therefore artifact-verifiable from its public repository and
fully replayable only with the specified captures. A newer download covering
the same dates is not an interchangeable raw input: packet order, snapshot
timing, and capture gaps are part of the recorded experiment.

\subsection{Install and verify mechanics}

Python 3.11 is the supported reference environment. From a complete Git clone:

\begin{lstlisting}[language=bash]
python3.11 -m venv .venv
source .venv/bin/activate
python -m pip install -r requirements.txt
python -m pip install --no-deps --editable ".[analysis,dev]"
env PYTHONPATH=. pytest -q
env PYTHONPATH=. python scripts/demo_event_driven_execution.py
env PYTHONPATH=. python scripts/verify_v2_artifacts.py
\end{lstlisting}

Native builds require a C++17 compiler, CMake 3.18 or newer, OpenSSL development
headers/libraries, and CPython development support for the selected
interpreter. The Make target builds the native tests, CLI, and Python binding:

\begin{lstlisting}[language=bash]
make cpp-test PYTHON=python
\end{lstlisting}

For Conda-provided OpenSSL, pass
\code{CMAKE_ARGS="-DOPENSSL_ROOT_DIR=$CONDA_PREFIX"}. The extension is built
for that Python interpreter and is not assumed portable across Python ABIs.
The adapter loads the build-directory extension, or the exact module selected
by \code{L2MM_CPP_MODULE}. \code{L2MM_CPP_BINARY} similarly selects the CLI
for standalone parity tests. Explicit selection fails if the requested
artifact is absent or invalid.

\subsection{Verify V3 without raw captures}

The recorded source commit must be available in Git history; shallow clones
must first obtain that history. A later tooling commit is expected to differ
from the research source, so the historical-source choice is explicit:

\begin{lstlisting}[language=bash]
env PYTHONPATH=. python scripts/verify_v3_artifacts.py \
  --source-revision recorded --artifacts-only
env PYTHONPATH=. python scripts/summarize_v3_development.py \
  --source-revision recorded --artifacts-only
\end{lstlisting}

Omit \code{--artifacts-only} when the selected raw captures are present to
hash-check those bytes as well. Omit \code{--source-revision} only when the
current working Python source is intended to be exactly the recorded research
source. The verifier never rewrites the run manifest to accommodate later
code changes.

\subsection{Reproduce the frozen research run}

The original V3 experiment is identified by the full commit below. An isolated
checkout avoids mixing later backend or tooling changes with the frozen
experiment:

\begin{lstlisting}[language=bash]
git worktree add --detach ../l2mm-v3-research \
  04df6294872a256374f910d12d5e5e3e08e8757a
\end{lstlisting}

Restore the selected captures under that checkout's \code{data/raw/} layout
using \code{DATA.md} and the inventory. With the project environment active,
the original serial, resumable entry point is:

\begin{lstlisting}[language=bash]
env PYTHONPATH=. python scripts/run_l2_panel.py --phase all
\end{lstlisting}

The completed research run scheduled independent commands concurrently through
the same runner checks; OFI commands sharing a cache remained sequential.
Scheduling did not change replay parameters or algorithms. The durable
manifest contains all 61 completed steps and 330 output identities. Ephemeral
caches and progress files are not substituted for durable evidence.

\subsection{Use and validate the native backend}

The current checkout exposes \code{--book-backend cpp} on replay and analysis
entry points. Native output and cache identities remain separate from the
frozen Python research result. Command help describes the input-specific
arguments for each entry point:

\begin{lstlisting}[language=bash]
env PYTHONPATH=. python scripts/run_replay.py --help
env PYTHONPATH=. python scripts/run_l2_panel.py --help
\end{lstlisting}

Full development pipeline parity compares both queue endpoints on the selected
120 hourly inputs. Worker count controls independent validation jobs, not
concurrent timing repetitions:

\begin{lstlisting}[language=bash]
env PYTHONPATH=. python scripts/validate_native_pipeline.py \
  --workers 2 \
  --output results/native_pipeline/development_parity.json
\end{lstlisting}

The recorded-hour application measurement is separate from standalone book
timing and reruns stream parity before accepting durations:

\begin{lstlisting}[language=bash]
env PYTHONPATH=. python scripts/validate_native_pipeline.py \
  --starts 2026-04-12T09:00:00 --benchmark-repeats 3 \
  --output results/native_pipeline/development_hour_benchmark.json

make cpp-benchmark PYTHON=python
\end{lstlisting}

These commands remeasure the local system and may produce different durations
from the committed benchmark. Correctness identities are expected to match;
wall-clock timings are observations of a particular environment and load.

\subsection{Build the complete report}

Tectonic 0.17.0 is the pinned TeX engine. The generated figures and numeric
macros use committed research and benchmark artifacts. The complete report's
targets preserve the historical V2 PDF:

\begin{lstlisting}[language=bash]
make complete-report
make complete-report-check
\end{lstlisting}

The output is \code{report/l2_mm_system_complete_report.pdf}; its source is
\code{report/v3/main.tex}. The existing \code{make report} builds the
historical \code{report/l2_mm_research_report.pdf}. Report checks inspect the
generated assets, PDF structure and text, embedded fonts, build diagnostics,
and recorded output identity. Poppler supplies \code{pdfinfo},
\code{pdffonts}, and \code{pdftotext}; \code{qpdf} can add a structural check.

\subsection{Canonical identities}

\begin{longtable}{@{}L{0.28\textwidth}L{0.67\textwidth}@{}}
\caption{Distinct hashes identify protocol, selection, source, and durable outputs.}\\
\toprule
Object & Identity \\
\midrule
\endfirsthead
\toprule
Object & Identity \\
\midrule
\endhead
\bottomrule
\endfoot
Integrity inventory logical identity & \nolinkurl{a3a99b0a616abe3bc39e0863ed047f075db9ed5d8118ced57c16140b499d8a61} \\
Selected-panel identity & \nolinkurl{760c55b7c0929b4a99657f6ca02eb723d930b9f48ebd3786d57bcb1a0f481122} \\
Development CSV SHA-256 & \nolinkurl{c779138fdffb739715c53cfa27c75b3c8ca140bc4f5a6128f60f2251948bd362} \\
V3 protocol SHA-256 & \nolinkurl{e1890d2e4fa715154f28778febe3367d391a8c8be1b44b766b4ccc0decf3e395} \\
V3 research Git commit & \nolinkurl{04df6294872a256374f910d12d5e5e3e08e8757a} \\
V3 Python source fingerprint & \nolinkurl{45141ba610bb070217e6d08bbb5fef9ca1a725f5c7d5e40b6e14881450a4cb4d} \\
V3 durable manifest SHA-256 & \nolinkurl{d70195704c98a14f7ce0daf747b97906754d7d2ffebcc181d308dad8af57d27a} \\
\end{longtable}

The inventory identity is a logical content identity, not the SHA-256 of its
JSON file. The source fingerprint hashes the selected Python source tree
rather than the compiled extension. Native evidence separately records its
source and binary identities. Later integration code does not replace
\ResearchCommit{} as the source of the V3 economics.

\subsection{Artifact map}

\begin{longtable}{@{}L{0.47\textwidth}L{0.48\textwidth}@{}}
\caption{Primary files for independent inspection. Paths are relative to the repository.}\\
\toprule
Path & Purpose \\
\midrule
\endfirsthead
\toprule
Path & Purpose \\
\midrule
\endhead
\bottomrule
\endfoot
\path{notebooks/v3_development_protocol.md} & Prospective hypotheses, frozen parameters, primary gate, bootstrap procedure, and stopping boundary. \\
\path{results/panels/btcusdt_l2_panel_v2/} & Frozen inventory, development/holdout selection, and historical V2 empirical evidence. \\
\path{results/panels/btcusdt_l2_panel_v3_event_driven/ARTIFACT_MANIFEST.json} & Complete V3 workflow, source and raw identities, and durable artifact hashes. \\
\path{results/panels/btcusdt_l2_panel_v3_development_summary/summary.json} & Recomputed decision, both endpoint diagnostics, uncertainty, and descriptive legacy comparison. \\
\path{results/replay_correctness/snapshot_timing_panel24.json} & Selected-depth snapshot timing audit and observed legacy exposure. \\
\path{cpp/parity_contract.md} & Frozen native input, operation, serialization, and arithmetic scope. \\
\path{results/cpp_orderbook/development_hour_parity.json} & Every-state recorded depth parity, without performance timing. \\
\path{results/cpp_orderbook/synthetic_20000.json} & Corrected standalone synthetic book benchmark. \\
\path{results/cpp_orderbook/development_hour.json} & Corrected standalone recorded-hour book benchmark. \\
\path{results/native_pipeline/development_parity.json} & Full replay and diagnostic stream comparison at the native backend boundary. \\
\path{results/native_pipeline/development_hour_benchmark.json} & Separately scoped application-path timing with parity gates. \\
\path{report/v3/generated/asset_manifest.json} & Report presentation inputs and generated asset identities. \\
\bottomrule
\end{longtable}

The root \code{results/} tree also includes exploratory and superseded
outputs. Use \code{results/README.md}, the versioned protocol, and the
canonical manifest before citing a file. An attractive number elsewhere in
the tree is not automatically part of the completed study.
